% Fragment intended to be \input into a larger document, alongside
% tool-modifications.tex. It assumes the preamble provides:
%
%   \usepackage{todonotes}   % for \todo[inline]{...}
%   \usepackage{booktabs}    % for \toprule / \midrule / \bottomrule
%
% Everything else relies only on core LaTeX: sectioning, \texttt, \emph,
% itemize, tabular and verbatim. Adjust the sectioning level to taste.

\section{The Evaluation Platform}
\label{sec:evaluation-platform}

The preceding section described the code generation solution under evaluation
and the two interventions it required. This section describes the apparatus
built around it: what the platform has to do, why it is assembled from an
existing orchestrator rather than written from scratch, and what each of the
decisions taken along the way was decided against.

\subsection{What the Platform Is For}
\label{sec:platform-purpose}

The research question is comparative. We want to say how a given code
generation solution performs on a given benchmark, and to be able to repeat
that statement for a second solution, a second model, or a second benchmark
without the comparison losing its meaning. Everything the platform does follows
from that, and can be stated as four obligations.

\begin{itemize}
\item \emph{Isolation.} A generated project is untrusted code that must be
      executed to be graded. It runs in a container, and so does its grader.
\item \emph{Fidelity.} The benchmark grades the result, not us. Whatever
      scoring procedure the benchmark's authors published is the procedure that
      runs, unmodified, on the artefact the solution produced.
\item \emph{Substitutability.} The code generation solution is the independent
      variable. Replacing it must not require touching the benchmark, the
      grading path, or the record-keeping.
\item \emph{Provenance.} Every recorded number must be traceable to the exact
      setup that produced it: benchmark version, task selection, tool revision,
      model, and hyperparameters.
\end{itemize}

\todo[inline]{State the research question in your own terms here, and say
explicitly which variables the thesis intends to vary (solution, model,
benchmark, hyperparameters) and which are held fixed. The four obligations
above are derived from that framing, so it should precede them.}

The scope is deliberately narrow. The platform is not a leaderboard, a service,
or a general agent framework; it is a reproducible experiment runner for a
thesis, and every feature below exists because one of the four obligations
demanded it.

\subsection{Why Harbor}

The first and largest decision was to not build an orchestrator. Containerised
benchmark execution is a solved problem with several mature implementations,
and reimplementing it would have consumed most of the available time while
producing something less trustworthy than what already exists. The question was
therefore which existing system to adopt.

The candidates fall into three families. There are \emph{benchmark-specific
harnesses}, such as the evaluation script shipped with SWE-bench or the
generation-and-evaluation pipeline shipped with NL2RepoBench itself: these run
one benchmark extremely well and nothing else. There are \emph{evaluation
frameworks} such as Inspect AI, which are general over tasks, models and
sandboxes, but treat the benchmark as something the experimenter expresses in
the framework's own vocabulary rather than something obtained ready-made. And
there are \emph{agent benchmark orchestrators} such
as Harbor, which take the container, the agent, and the grader as the three
things they schedule.

We assessed the candidates against six criteria, weighted by how directly each
serves the four obligations of Section~\ref{sec:platform-purpose}.

\begin{table}[htbp]
\centering
\caption{Decision matrix for the orchestrator. Scores run from 0 (does not
provide) to 3 (provides directly), weighted by the criterion's importance to
the obligations in Section~\ref{sec:platform-purpose}. The benchmark-native
harnesses --- SWE-bench's evaluation script and NL2RepoBench's own pipeline ---
are scored as one class, since they differ in maturity rather than in kind:
each grades one benchmark well and offers no route to a second.}
\label{tab:orchestrator-matrix}
\begin{tabular}{lcccc}
\toprule
                                        & \multicolumn{4}{c}{Candidate} \\
Criterion (weight)                      & Harbor & Bespoke & Benchmark- & Inspect \\
                                        &        &         & native     & AI      \\
\midrule
Containerised, isolated execution (3)   & 3  & 2  & 3  & 3  \\
Agent-agnostic adapter interface (3)    & 3  & 3  & 1  & 2  \\
Benchmark-agnostic, many benchmarks (3) & 3  & 1  & 0  & 2  \\
Version pinning and provenance (2)      & 3  & 2  & 1  & 2  \\
Concurrency and scale (2)               & 3  & 1  & 3  & 3  \\
Maintenance burden on this work (2)     & 3  & 0  & 2  & 1  \\
\midrule
Weighted total (maximum 45)             & 45 & 24 & 24 & 33 \\
\bottomrule
\end{tabular}
\end{table}

Two of the columns need their scores justified, since neither is obvious.

The \emph{bespoke} column scores full marks exactly where the system would be
built to fit --- the agent interface --- and poorly everywhere the fit is
irrelevant. The scale of what would be reimplemented is worth stating
concretely: the adopted release carries some thirty-three thousand lines across
forty-two environment-provider modules, covering Docker and Compose alongside
Kubernetes, EC2, Modal and roughly two dozen other backends, together with a
typed retry policy, per-phase timeout multipliers, a task registry, a results
database and a viewer. A bespoke runner would reproduce a thin slice of this,
and the slice it reproduced would be the one we had happened to need so far. It
scores 2 rather than 0 on isolation because shelling out to Docker Compose is
genuinely within reach --- the platform does exactly that for image
preparation --- and 2 rather than 3 on provenance because run-level records are
easy to design well, while the content-addressed benchmark registry that makes
version pinning meaningful is not something we would have built.

The \emph{Inspect AI} column is the closest competitor, and its initial
assessment here was wrong in a way worth correcting rather than quietly fixing.
Inspect is not, as we first assumed, restricted to prompt-and-score
interactions: it provides first-class sandboxing, with per-sample Docker
Compose configurations, network isolation by default, and external providers
for Kubernetes and several cloud backends. It scores 3 on isolation and 3 on
concurrency accordingly. Its provenance is also stronger than a framework of
its kind usually offers --- an evaluation log records the git origin and commit
of the evaluation code and the versions of installed packages --- but the
dataset it records is identified by location and sample count rather than by
content digest, which is precisely the guarantee the fidelity obligation leans
on, so it scores 2 rather than 3.

What decides against Inspect is the pair of remaining criteria. Its collection
of packaged evaluations is large and includes SWE-bench, but not NL2RepoBench,
so the benchmark would have to be ported into Inspect's task-and-scorer form
--- 104 tasks, each with a compose file and a tester sidecar --- which is the
integration effort the dependency model was adopted to avoid. And its mechanism
for hosting a third-party agent, the agent bridge, works by redirecting the
agent's own model calls through Inspect's provider; by default that path
discards generation parameters such as \texttt{max\_tokens} and
\texttt{temperature} on the grounds that they were tuned for the agent's native
model. Those are among the hyperparameters this work varies and records, so the
bridge would have to be worked against rather than with.

\todo[inline]{If you did not in fact evaluate Inspect AI at the time, soften
the framing above from ``our initial assessment was wrong'' to a plain
description of the alternative. The correction reads well only if the
assessment actually happened.}

\todo[inline]{Confirm the candidate list. These are the alternatives that are
defensible in writing, but you know which ones were genuinely on the table
during the project. Remove any that were never considered, and add any that
were --- an internal tool, a supervisor's suggestion, or a system you tried and
abandoned carries more weight in a thesis than a hypothetical.}

\todo[inline]{The weights and scores are a defensible reading of the evidence,
not a measurement. Either adjust them to match your own judgement, or add a
sentence conceding that the matrix formalises a qualitative judgement rather
than replacing it. Reviewers are sceptical of decision matrices whose totals
happen to confirm the choice already made; saying so first disarms the
objection.}

Three of the six criteria decided it in practice.

\emph{Agent-agnosticism} is the one the thesis cannot do without. Harbor treats
the agent as a plugin: a class with an \texttt{install} step and a \texttt{run}
step, given a container and an instruction, judged solely by the state of the
workspace it leaves behind. This is exactly the shape of the substitutability
obligation. A benchmark's own harness, by contrast, has an agent baked in ---
NL2RepoBench generates with OpenHands --- and evaluating a different solution
through it means either replacing that path or reproducing the grading
elsewhere. The former is invasive, the latter forfeits fidelity.

\emph{Benchmark-agnosticism} decides the second axis of comparison. Harbor
carries whole benchmarks in a registry as datasets pinned by content digest, so
a second benchmark is a change of one line in a configuration file rather than
an integration project. Committing to one benchmark's harness would fix the
thesis to one benchmark permanently.

\emph{Maintenance burden} decides against building it ourselves. The bespoke
option scores well on adapter design, for the trivial reason that we would
design it, and mid-table at best on everything that has to exist regardless of
who the experiment is for: container lifecycle, image handling,
concurrency, retry, result storage and a results viewer are all work that
produces no research finding. The project brief for this work states the
position directly --- Harbor is the orchestrator, and a custom platform is out
of scope.

The cost of the decision is a dependency on Harbor's model of a task and on its
release cadence. We accept it and mitigate it by pinning: the project depends
on \texttt{harbor==0.21.0} exactly, so an upstream release cannot silently
change what a run does.

\todo[inline]{If Harbor was a given rather than a choice --- prescribed by a
supervisor, a chair, or an existing group setup --- say so plainly here instead
of presenting the matrix as an open evaluation. A retrospective justification
of a constraint is legitimate and worth writing, but it should be labelled as
one.}

\subsection{What the Platform Adds}

Adopting an orchestrator leaves a specific and much smaller remainder: the
things Harbor does not do because they are particular to this experiment. The
platform is four modules and a launcher, and each exists for one of them.

\begin{itemize}
\item \texttt{experiment\_config.py} --- the experiment setup as a validated,
      archivable object.
\item \texttt{benchmark.py} --- resolving a benchmark to the exact tasks a run
      will execute, and making their images obtainable.
\item \texttt{self\_collaboration\_agent.py} --- the Harbor agent adapter,
      running outside the container.
\item \texttt{run\_self\_collaboration.py} --- the runner, uploaded into the
      container and executed there.
\item \texttt{main.py} --- the launcher that composes the three and hands
      Harbor a job.
\end{itemize}

The split between the two agent-side files deserves a note, since it is not
obvious from the file names. The adapter runs on the host, in Harbor's process,
and can therefore see the experiment configuration but not the workspace; the
runner runs inside the task container, and can see the workspace and the code
generation tool but nothing of Harbor. Everything crossing that boundary
crosses as a file: the instruction and the hyperparameters inward, the logs and
the usage record outward. Keeping the boundary explicit is what makes the
in-container half testable without Docker.

\subsection{Configuration as the Unit of an Experiment}

\subsubsection{One file describes one experiment}

The rule the platform enforces is that everything which changes what a run does
lives in one file under \texttt{configs/}, and \texttt{.env} holds credentials
and nothing else. A configuration names the benchmark and the task selection,
the model backend, the revision of the code generation tool, and the
hyperparameters its three roles receive.

The alternative --- command-line flags, environment variables, or defaults
buried in the adapter --- was rejected for a reason specific to the provenance
obligation. A setup spread across three mechanisms cannot be archived, and a
result whose setup cannot be archived cannot be compared to anything later. By
putting the setup in one document, the launcher can write the resolved form of
that document next to the job before the run begins, and any recorded number
remains traceable to the file that produced it.

This also gives the comparison protocol a mechanical form. To evaluate a second
code generation solution, one copies a configuration, changes only the
\texttt{agent} section, and leaves the \texttt{benchmark} block
byte-identical. That the two runs saw the same benchmark is then a property one
can verify by inspection rather than a claim one has to trust.

\subsubsection{Validation before Docker starts}

Configurations are validated strictly and early: unknown keys in any section
are rejected rather than ignored, hyperparameters are type-checked against a
declared schema, an unset credential is an error, and a configuration must
choose exactly one of a registry benchmark and a local task directory.

Strictness here is a deliberate trade against convenience. A typo in a
hyperparameter name --- \texttt{max\_round} for \texttt{max\_rounds} --- is,
under a permissive reader, an experiment that silently runs with the default
and reports a number attributed to a setting that never took effect. That
failure is invisible in the results and expensive to detect afterwards.
Rejecting the key costs one error message and removes the failure mode
entirely. The same reasoning covers the endpoint check: DeepSeek's
OpenAI-compatible API has no \texttt{/api/v1} path, and a run pointed at one
fails several minutes and several container pulls later, with an error that
does not name the cause.

The template syntax \texttt{\$\{VAR\}} and \texttt{\$\{VAR:-default\}} exists
so that the model backend can be overridden from the environment without
editing a tracked file. The archived snapshot stores the \emph{resolved} value,
never the template, so the record says which model answered rather than which
environment variable was consulted.

\subsubsection{Hyperparameters as a declared schema}

The hyperparameters of the three roles --- round count, per-role step budgets,
sampling parameters, and the Tester's test command --- are declared centrally
with their types, and defaults are filled in at load time so that what is
archived is the \emph{effective} setup rather than only the part that was
written down.

One hyperparameter deliberately has no default. Omitting \texttt{test\_command}
disables the Tester phase entirely, which, as the previous section explained,
reduces the published three-role method to a two-role one. That is a legitimate
ablation, but it is not something a configuration should fall into by leaving a
key out, so the platform requires it to be stated either way.

\todo[inline]{Say which of these hyperparameters you intend to treat as
experimental variables and which are fixed for all runs. The platform makes all
of them varied-and-recorded; the thesis presumably varies far fewer, and the
values it fixes (\texttt{max\_rounds: 3}, \texttt{temperature: 0.0}, the step
budgets) need a justification here --- particularly the temperature, since
determinism is a methodological choice with consequences for how variance is
reported.}

\subsection{The Benchmark as a Dependency}

\subsubsection{Not checked in}

A benchmark enters the repository as a pinned dependency rather than as copied
task files. Harbor's registry carries NL2RepoBench as a digest-pinned dataset
of 104 tasks; a configuration names the dataset, the version, and which tasks
to run, and Harbor downloads the rest.

Vendoring the tasks was the obvious alternative and is worse in three ways. It
makes the repository the authority on what the benchmark contains, so a local
edit --- however well intentioned --- silently changes the benchmark under
which results were reported. It obscures the version, since a copy has no
digest. And it makes scaling from one task to the full benchmark an integration
effort rather than what it should be, a filter in the configuration. Under the
dependency model, the difference between the smoke-test configuration and the
full run is the presence of one \texttt{task\_names} entry.

The pinning is done twice over. The configuration carries a content digest,
and if a run is given a floating reference instead, the launcher resolves it
once at launch and re-pins the result before any trial starts. Every trial of a
run therefore sees one version of the benchmark, even if the registry moves
mid-run.

\subsubsection{Selection performed by Harbor, not beside it}

The preparation step has to know which tasks a run will execute, in order to
prepare their images. It could determine this by reading the configuration and
applying the filters itself. It does not: it constructs Harbor's own dataset
configuration object and asks it, so that the selection performed during
preparation is the same code path as the selection performed by the job.

This closes a class of bug that would otherwise be nearly undetectable. Two
independent implementations of glob matching, exclusion, and truncation agree
on every case one thinks to test and disagree on some case one does not, and
the symptom --- a task whose image was never prepared --- surfaces as a
container failure deep into a long run. Delegating removes the possibility of
disagreement rather than reducing it.

\subsubsection{The image mirror}

Most Harbor benchmarks need nothing beyond the above. NL2RepoBench is the
exception, and the exception is instructive because of how it was handled.

Every one of its 104 tasks names its tester image on a private registry mirror
requiring credentials this project does not have. The identical images are
published openly under the project's own namespace on GHCR. The benchmark, as
distributed, therefore cannot be run by anyone outside the organisation that
published it, despite all of its parts being public.

Three remedies were available. We could patch each task's compose file to name
the reachable image, which modifies the benchmark and thereby forfeits the
fidelity obligation --- and does so 104 times. We could exclude NL2RepoBench,
which forfeits the benchmark. Or we could satisfy the reference rather than
change it. The platform does the third: before a run, it resolves the tasks the
job will execute, reads the image references they name, pulls each from its
public home, and tags it locally under the name the task expects. Docker
Compose uses a locally present image without contacting a registry, so the
benchmark runs unmodified and is never aware that anything happened.

The rule is expressed as a prefix rewrite in the configuration --- one rule
covers all 104 tasks --- and is optional and benchmark-specific by design. A
benchmark whose images are reachable declares no rule, and the preparation step
reduces to resolving and pinning. The digest of every image actually used is
recorded with the run, so the substitution is documented rather than merely
performed.

\todo[inline]{Consider whether to report this upstream, and say here whether
you did. A benchmark that cannot be executed outside its authoring organisation
is a reproducibility defect in a published artefact, and noting it --- with or
without a filed issue --- is a small independent contribution of this work.}

\subsection{Running the Solution Inside the Task}

\subsubsection{Cloning the tool at a pinned commit}

The agent adapter installs the code generation solution by cloning it into the
container at the exact commit named in the configuration, rather than mounting
the working copy checked out beside the repository. The submodule under
\texttt{code\_generation/} exists for reading and reference; no run consumes
it.

This is a provenance decision. A mounted working copy makes the tool's revision
whatever the developer's checkout happened to be at launch --- unrecorded,
unreproducible, and quietly different between the run that produced a number
and the run meant to confirm it. Cloning a named commit makes the tool's
version part of the archived setup, on the same footing as the benchmark
digest. The runner then reports the commit the container actually ended up
with, which is discussed below.

\subsubsection{The instruction travels as a file}

The task instruction is written to a file and uploaded into the container
rather than passed as a command-line argument.

The immediate cause was encoding. NL2RepoBench instructions are long
natural-language specifications containing mathematical notation and non-ASCII
characters, and a command line is a lossy channel for them --- quoting, length
limits, and, on Windows hosts, a code page that is not UTF-8. The failure is
not clean: the instruction arrives corrupted rather than absent, and the
solution then works diligently on a subtly wrong specification. Writing a file
with an explicit encoding and reading it back with the same one removes the
channel from the problem. The hyperparameters travel the same way, for
consistency rather than necessity.

\subsubsection{Initialising the workspace repository}

The runner initialises a Git repository in the workspace, with one empty
commit, before the session starts.

This is the platform's half of the issue described in the previous section. The
code generation solution inspects the repository to decide whether its Coder
produced anything, and shows the Coder a diff of its own previous attempt
between rounds. Both facilities assume a repository exists, which is true of
the benchmark the solution was written against and false of a benchmark that
hands the agent an empty workspace. Rather than modify the solution to tolerate
a workspace without version control, the platform supplies the precondition the
solution already expects. The corresponding change on the solution's side ---
inspecting status rather than diff, so that untracked files are seen --- is the
subject of Section~\ref{sec:tool-modifications}, and the two are
complementary: the platform provides the repository, and the patch makes the
inspection correct within it.

\todo[inline]{Cross-reference: the \texttt{\textbackslash ref} above assumes
the tool-modifications section carries the label
\texttt{sec:tool-modifications}. Add that label to
\texttt{tool-modifications.tex}, or adjust this reference to whatever labelling
scheme the thesis uses.}

\subsubsection{Usage accounting and rate limits}

Cost and token consumption are part of what a comparison between solutions
should report, and the solution does not record them. The runner obtains them
by wrapping the single function through which the solution reaches the model,
accumulating the usage figures each response carries, and writing the totals
where Harbor reads them back into its results view.

Wrapping rather than editing was chosen for the same reason the previous
section gives for keeping modifications minimal: the fewer edits the tool under
measurement carries, the more credible the claim that what is measured is the
tool. The wrapper observes; it does not alter the request, the response, or
anything the solution does with either. It is also the natural place for the
retry policy, since a rate-limited request is a property of the endpoint rather
than of the method. The solution exhausts its own three attempts, the wrapper
allows two further rounds with increasing backoff, and if all nine requests are
throttled the run fails with a message that names the rate limit and the model
rather than a generic failure to call an API.

Cost is recorded only when the endpoint reports it in its usage response. Where
it does not, the field is left empty rather than estimated from a local pricing
table, on the grounds that an estimate drawn from a table that ages silently is
worse than an acknowledged gap.

\todo[inline]{Decide how cost enters the evaluation, if at all. If the thesis
compares solutions on cost as well as on reward, the empty-when-unreported
policy constrains which providers can be used for those runs, and that
constraint belongs in the methodology.}

\subsection{Provenance and Results}

Harbor writes each job beneath \texttt{jobs/}, and the platform adds records of
its own so that a job directory is self-describing. Three of them carry the
argument.

The \emph{archived configuration} is the resolved experiment setup, written
before the run starts, with templates expanded and defaults filled in. Writing
it up front rather than on completion is deliberate: a crashed run is precisely
the case where one needs to know what was attempted. It names the credential's
environment variable and never its value.

The \emph{benchmark record} states the version the benchmark resolved to, the
tasks the run selected, and the content digest of every tester image used. It
is the evidence for the fidelity claim: it says which grader ran, not merely
which grader was intended.

The \emph{resolved setup}, written from inside the container, states what
actually ran --- the commit of the code generation tool the container ended up
with, the model the environment resolved to, the effective hyperparameters, and
whether the Tester was enabled. Its purpose is to be redundant with the
archived configuration. The configuration records intent and the resolved setup
records observation, and a discrepancy between them is a bug the platform
surfaces rather than conceals.

Alongside these, each trial retains the generated workspace, the solution's
console log and structured session history including each round's test result,
and the benchmark's own verifier output with the reward. The reward is the
fraction of the task's hidden reference tests that passed.

\todo[inline]{State the metric the thesis reports and how it aggregates across
the 104 tasks --- mean pass fraction, resolve rate at a threshold, or both ---
and whether a single attempt per task is sufficient given a temperature of
zero. If multiple attempts are used, say how they are combined.}

\subsection{How a Task Is Graded}

Each task runs two containers. In \texttt{main}, a generic Python and Node
image, the solution generates the project under the workspace. Beside it runs a
\texttt{tester} sidecar built from the benchmark's own evaluator image, holding
the hidden reference tests. When the solution finishes, Harbor's verifier hook
signals the sidecar over the shared workspace volume; the sidecar strips any
test files the agent generated, copies the remaining code on top of its own
reference tests, installs the package, runs the suite, and reports the passing
fraction as the reward.

Two properties of this arrangement are load-bearing. The reference tests are
never present in the container the agent works in, so they cannot be read,
imported, or accidentally satisfied by being overwritten --- and the sidecar's
discarding of agent-authored test files closes the remaining route, in which a
generated test file shadows a reference one by name. And the procedure is the
benchmark's own upstream flow rather than a reimplementation of it, which is
what the fidelity obligation requires: a reimplementation would be a second
grader whose agreement with the first is an assumption.

The Tester role inside the solution runs an entirely separate command, in the
agent's own workspace, over the code and the tests the agent itself wrote. It
is self-verification, and it has no access to and no knowledge of the grading
path described here.

\subsection{Testing the Platform}

The platform carries a test suite of its own, which is worth a sentence because
of what it covers. The tests do not exercise the code generation solution or
the benchmark; both are dependencies with their own correctness. They cover the
seams: that hyperparameters written in a configuration arrive as the values
Harbor records and the runner receives, that an unknown key is rejected rather
than dropped, that the archived setup contains no credential, that the image
mirror pulls only images that are genuinely absent and leaves present ones
alone, that a floating benchmark reference is pinned to what it resolved to,
and that the runner omits the Tester exactly when no test command is
configured.

This is the appropriate target. Every one of these is a failure that would
otherwise produce a plausible-looking number attributed to a setup that did not
run, which is the one category of defect a comparative evaluation cannot
tolerate.

\subsection{Limitations}

Three limitations follow directly from the decisions above, and are better
stated than discovered.

The platform inherits Harbor's model of a task. A benchmark that does not fit
that model --- interactive tasks, tasks without a container, tasks whose
grading is not the exit status of a command --- is out of reach without exactly
the work the adoption decision was taken to avoid.

The image mirror is a workaround for a defect in one benchmark's distribution.
It is small, declarative, and recorded, but it remains a rule asserting that
two image references denote the same image, and that assertion is checked by
digest at run time rather than guaranteed in advance.

Finally, the fidelity obligation is satisfied with respect to the benchmark's
grading procedure, not with respect to the benchmark's intended generation
procedure. NL2RepoBench's published results were produced by a different agent
under that benchmark's own harness; ours are produced by a different solution
under Harbor. The comparison the platform supports well is between solutions
run through this platform. Comparison against externally published figures is a
weaker claim, and should be made carefully if at all.

\todo[inline]{This last limitation is the one most likely to be challenged in a
defence. Decide the position now: either the thesis compares only solutions run
through this platform --- in which case say so in the methodology and drop any
table that places your numbers beside published ones --- or it makes the
external comparison and needs a paragraph on why the difference in harness does
not account for the gap.}

\todo[inline]{Add any limitation arising from the scale actually achieved. If
the full 104-task run was not completed, or was completed only for a subset of
models, that constrains every claim in the results chapter and belongs here
rather than in a footnote.}
